\chapter{Background}\label{ch:literature-review}

This chapter gives the background the argument uses, in dependency
order: evaluation benchmarks and arena platforms
(Section~\ref{sec:benchmarks}); the family the construction belongs to,
stated as positioning rather than surveyed
(Section~\ref{sec:taxonomy-construction}); what is known to go wrong
with LLM judges (Section~\ref{sec:llm-as-a-judge}); and pairwise
ranking models (Section~\ref{sec:ranking-models}).
Section~\ref{sec:routing} closes with query routing, the downstream
task for which per-domain capability profiles are needed, and states
the narrowed gap this thesis defends against the routing systems that
already induce structure of their own.

The sections are not sized by the size of their literatures. They are
sized by how much of each this thesis's argument consumes.
Section~\ref{sec:llm-as-a-judge} is the longest because
Chapter~\ref{ch:results}'s central result is a measurement of a judge,
and Section~\ref{sec:taxonomy-construction} is the shortest because the
construction is an instrument.




\section{LLM Evaluation Benchmarks}\label{sec:benchmarks}

\subsection*{Static Benchmarks and Their Limitations}

The first generation of Large Language Model (LLM) evaluation relied on
static instruments such as \textit{Massive Multitask Language Understanding}
(MMLU)~\parencite{hendrycks2021measuring} and the \textit{Holistic Evaluation
of Language Models} (HELM) framework~\parencite{liang2022holistic}.
MMLU uses over \num{15000} multiple-choice questions across 57 subjects to
probe professional and academic competence; at release, the best-performing
model (GPT-3, 175B) reached 43.9\% accuracy against an estimated human
expert ceiling of 89.8\%~\parencite{hendrycks2021measuring}.
HELM attempted a more multidimensional approach by evaluating 30 models
across 42 scenarios and seven metrics (accuracy, calibration,
robustness, fairness, bias, toxicity, and efficiency), making it the
first large-scale effort to resist the reduction of model quality to a
single number~\parencite{liang2022holistic}.

In both frameworks the evaluation corpus and its categorical
organisation are fixed at construction time.
This creates three compounding failure modes.
First, \textit{saturation}: as model capabilities advance, score
distributions compress toward the ceiling and the benchmark loses
discriminative power. Frontier models now score at or near MMLU's
estimated human-expert ceiling of 89.8\%, which leaves the benchmark
unable to rank them against each
other~\parencite{hendrycks2021measuring, wang2024mmlupro}.
Second, \textit{contamination}: benchmark questions circulate online and
enter training corpora; researchers found that ChatGPT and GPT-4 could guess
missing answer options in MMLU at rates of 52\% and 57\% respectively,
far above what genuine reasoning would predict~\parencite{deng2024contamination}.
Third, \textit{categorical brittleness}: editorially defined subject
categories cannot adapt to the capability space of new model generations,
leaving blind spots at boundaries and failing to capture cross-domain
dimensions that emerge only after the benchmark is fixed.
HELM's fixed scenario taxonomy, despite its breadth, inherits the same
limitation: its 42 scenarios were defined for the model generation of 2022
and do not evolve with the capability frontier~\parencite{liang2022holistic}.

\subsection*{From MMLU to MMLU-Pro}

The saturation and contamination of MMLU motivated a line of progressively
harder benchmarks.
\textit{BIG-Bench Hard} (BBH)~\parencite{suzgun2022bigbench} addressed
discriminability by extracting the 23 BIG-Bench tasks on which prior models
had not outperformed average human raters, covering multi-step arithmetic,
logical deduction, causal judgment, and temporal reasoning across 6,511
examples.
BBH demonstrated that chain-of-thought (CoT) prompting enabled
PaLM to surpass human-rater performance on 10 of those 23 tasks, a result
that confirmed both the diagnostic value of the harder subset and the
inadequacy of the four-choice format for measuring genuine
reasoning~\parencite{suzgun2022bigbench}.

\textit{MMLU-Pro}~\parencite{wang2024mmlupro} extends this logic directly to
the MMLU lineage.
It expands the answer set from four to ten options per question,
eliminates trivial and noisy items from the original MMLU, and integrates
reasoning-focused questions drawn from academic exams and textbooks across
14 domains (mathematics, physics, chemistry, engineering, law, and
computer science among them), yielding \num{12000} rigorously curated
items~\parencite{wang2024mmlupro}.
The effect on model scores is substantial: accuracy drops by 16\% to 33\%
relative to MMLU, and prompt sensitivity falls from 4--5\% to 2\%, indicating
that MMLU-Pro is more robust to surface variation and more discriminative of
underlying reasoning capability.

The problem of contamination, however, persists even in harder benchmarks
so long as questions are static.
\textit{LiveBench}~\parencite{white2024livebench} addresses this
structurally by refreshing its question pool on a rolling basis from
recently published sources, substantially reducing contamination risk
since questions are drawn from material published after common training
cutoffs.
This approach trades domain coverage for freshness, and its task set
remains narrower than MMLU-Pro's 14-domain structure.
Depth of domain coverage (MMLU-Pro) and resistance to contamination
(LiveBench) therefore pull against each other.
A framework that fixes neither the questions nor the domains, inducing
its evaluation structure from the query corpus instead, escapes
categorical brittleness, though its domain coverage remains bounded by
the diversity of the supplied corpus.


\subsection*{Pairwise Arena Evaluation Platforms}

A parallel line of work addressed a different limitation of static
benchmarks, the absence of human preference as an evaluation signal.
\textit{Chatbot Arena}~\parencite{chiang2024chatbotarena}, introduced by
Chiang et al.\ at UC Berkeley and LMSYS, is the canonical instance.
Users submit a free-form prompt, receive responses from two anonymous models
simultaneously, and vote for the preferred response; the platform has
collected over 240K votes across more than 50 models as of early
2024~\parencite{chiang2024chatbotarena}.
Pairwise comparisons are aggregated using a Bradley--Terry model, and
confidence intervals on the resulting coefficients are computed via the
sandwich robust standard error, providing statistically grounded rankings
with explicit uncertainty estimates.
To improve sample efficiency, Chatbot Arena employs an adaptive sampling
rule that allocates comparison budget proportionally to the reduction in
confidence interval size, reaching a given precision on the win matrix
with substantially fewer comparisons than uniform random
sampling~\parencite{chiang2024chatbotarena}.

The arena paradigm itself has since come under systematic critique.
\textcite{singh2025leaderboard} document that Chatbot Arena's operating
policies --- private testing of model variants before public listing,
asymmetric access to vote data, and unequal sampling and deprecation
rates across providers --- systematically advantage a small set of
participants, so that observed Arena rank is partly a function of
platform engagement rather than of capability alone.
The consequence relevant here is methodological: a leaderboard is a
measurement instrument, and its measurement properties must be audited
independently of the quantity it reports --- the stance this thesis
applies to its own evaluation chain.

Recognising that a single global ranking conceals domain-specific
capability variation, the Arena line of work has begun inducing
structure from its own query stream. The documented pipeline is
\textit{BenchBuilder}~\parencite{li2024benchbuilder}, which embeds
roughly 200K Arena prompts, clusters them hierarchically into some
4{,}000 topic clusters, scores prompts with an LLM against seven
quality criteria, and distils a 500-prompt benchmark
(Arena-Hard-Auto) whose model ranking reaches 98.6\% correlation with
the full human-vote leaderboard at a judging cost of roughly \$20 per
evaluated model; Arena's public
per-category leaderboards follow the same
embed--cluster--LLM-label recipe~\parencite{arena2025categories}.
BenchBuilder demonstrates that corpus-induced structure over an
evaluation query stream is practical at scale, which is precisely why
the gap this thesis addresses must be stated narrowly
(Sections~\ref{sec:taxonomy-construction} and~\ref{sec:routing}).

What the pipeline does not attempt is equally precise. Its clusters
are a means to prompt \emph{filtering}: they carry no probability
model, so cluster membership is hard, the coherence of a cluster is
never measured, and the hierarchy is discarded once the 500 prompts
are selected. There is no split-or-stop statistic that makes the
partition itself an auditable object, and no per-cluster model
statistics survive into the published benchmark.
TaxoArena differs on exactly this axis: its cells are vMF-modelled
regions of the embedding sphere whose coherence is statistically
tested, and each cell carries a Bradley--Terry profile with
Fisher-information confidence intervals --- the partition is the
measured object, not a preprocessing step.

\textit{Prompt-to-Leaderboard} (P2L)~\parencite{frick2025p2l} pushes
per-domain resolution to its continuous limit: a language model is
trained on Arena's human preference votes to map an individual prompt
directly to a vector of prompt-conditional Bradley--Terry
coefficients, yielding a leaderboard --- and hence a routing decision
--- per prompt; a P2L-based router reached first place on Chatbot
Arena in January 2025.
P2L is the strongest existing occupant of the per-domain
capability-profile territory, and Section~\ref{sec:routing} returns
to it as a competing baseline.
The differentiator is representational: P2L amortises the capability
profile into the weights of a regression model, so its ``domains''
are implicit, continuous, and inspectable only through the model's
outputs, whereas TaxoArena produces a discrete, human-interpretable
cell structure whose per-cell rankings are statistically
characterised and can be hypothesis-tested cell by cell.



\section{Taxonomy Induction: Positioning}\label{sec:taxonomy-construction}

The construction used in this thesis is an instrument rather than a
contribution (Section~\ref{sec:approach}), so this section says what
family it belongs to and what it borrows, and does not survey the
field. Three bodies of work are load-bearing for later chapters; the
rest is named and set aside.

\textbf{Set aside.}
Supervised hierarchical text classification (HiAGM,
HiTIN~\parencite{zhou2020hiagm, zhu2023hitin}) requires the hierarchy
as supervision, and the hierarchies it is trained against --- Reuters
editors for RCV1, domain librarians for WoS --- are editorial
artefacts whose granularity, balance and separability are not
reproducible. Such methods assume an already-coherent taxonomy and
learn to navigate it; they offer no construction method, and the two
supervised baselines originally planned for this thesis were removed as
unimplemented. Dasgupta's cost function~\parencite{dasgupta2016cost}
supplies a principled global objective for hierarchical clustering, but
it defines no local stopping rule, and this thesis claims no Dasgupta
optimality for its trees. Reference-dependent hierarchy metrics
($h\!F_\beta$ and its relatives~\parencite{kosmopoulos2015hmetrics})
require a gold hierarchy, and none exists for a corpus-induced taxonomy
over MMLU-Pro; they are not reported.

\textbf{Unsupervised induction, and the one differentiator that
matters.}
\textit{TaxoGen}~\parencite{zhang2018taxogen} builds a topic taxonomy
top-down by adaptive spherical clustering with local embedding
retraining at each node, and \textit{CoRel}~\parencite{huang2020corel}
extends the idea to the seed-guided setting by learning a
relation-transferring module that propagates a user-supplied skeleton
along embedding paths. TaxoArena shares CoRel's seed-guided starting
point and differs in kind: the MMLU-Pro labels initialise anchor
membership and are thereafter \emph{superseded} by geometric evidence
--- migration, splitting, and topology operations --- whereas CoRel's
relation-transferring module encodes the supplied skeleton as a
structural prior throughout construction. Both methods also operate at
the \textit{term} level, where a node is a cluster of words; here a
node must characterise the distributional geometry of full
natural-language queries, which needs a different node model, a
different splitting criterion, and a different coherence measure.

\textbf{The occupying systems, and the narrowed claim.}
The embed--cluster--LLM-label recipe that BenchBuilder applies to
benchmark curation (Section~\ref{sec:benchmarks}) has become an
industrial pattern for inducing taxonomies from live corpora, and any
claim that corpus-induced taxonomy for evaluation is unoccupied
territory would be false; this thesis does not make it.
\textit{TnT-LLM}~\parencite{wan2024tntllm} has an LLM generate and
iteratively refine a label taxonomy over conversation logs, then
distils it into lightweight classifiers deployed at scale over Bing
Copilot traffic. \textit{Clio}~\parencite{tamkin2024clio} embeds and
clusters millions of Claude conversations, has an LLM name each
cluster, and assembles the named clusters into a navigable hierarchy
for privacy-preserving usage monitoring.
\textit{TopicGPT}~\parencite{pham2024topicgpt} replaces the
clustering step itself with prompted topic generation and assignment.
All three, like BenchBuilder, induce human-readable structure from a
corpus; all three consume that structure for classification,
monitoring, or curation.
The claim this thesis defends is a conjunction that none of them
provides: TaxoArena combines (a) a generative node model --- vMF
distributions on the embedding sphere --- with a geometric
split-and-stop rule whose acceptance statistic is chance-corrected,
so that the coherence of every cell is a tested quantity rather than
an artefact of labelling, and (b) statistically characterised
per-cell capability profiles (Bradley--Terry coefficients with
Fisher-information confidence intervals) consumed by a router.
BenchBuilder, TnT-LLM, Clio, and P2L each occupy one half of this
conjunction --- taxonomy induction without statistical node
characterisation, or capability profiles without a discrete,
testable cell structure --- and none occupies both.

\subsection*{Spherical Embeddings and vMF Clustering}

When text is encoded by a neural model and $\ell_2$-normalised to the
unit hypersphere, cosine similarity becomes the natural inner product
and the von Mises--Fisher (vMF) distribution becomes the canonical
probability model for directional data~\parencite{banerjee2005vmf}.
The vMF density for a unit vector $\mathbf{x} \in \mathcal{S}^{d-1}$,
with mean direction $\boldsymbol{\mu}$ and concentration
$\kappa \geq 0$, is
\begin{equation}
    p(\mathbf{x};\boldsymbol{\mu},\kappa) =
    C_d(\kappa)\exp\!\bigl(\kappa\,\boldsymbol{\mu}^\top\mathbf{x}\bigr),
    \label{eq:vmf}
\end{equation}
where $C_d(\kappa)$ is a normalising constant involving a modified
Bessel function of the first kind.
Fitting a mixture of vMF distributions by EM is equivalent to
minimising a Bregman divergence on the sphere, and high-order
approximations to the concentration estimate remain numerically stable
in the high-$\kappa$, high-dimensional regime typical of neural
embedding spaces~\parencite{banerjee2005vmf, hornik2014vmf}.
This is the node model used throughout; the estimation details,
including the small-sample stabilisation of $\kappa$ and the reason
cross-node comparisons must use a \emph{shared} concentration, are in
Appendix~\ref{app:dag-formalism}.

The recursive construction additionally needs an embedding space in
which coarse domain-level and fine sub-topic distinctions are both
resolvable within one representation. Matryoshka Representation
Learning~\parencite{kusupati2022mrl} provides it by optimising a single
embedding jointly across nested prefix dimensions, so the most
discriminative directions concentrate in the leading coordinates and
any prefix is a coherent representation rather than a truncation
artefact. TaxoArena uses a single fixed 256-dimensional slice at every
depth rather than a per-depth schedule, which places every node on the
same sphere and makes cross-depth comparisons directly valid; the
justification is developed in Appendix~\ref{app:dag-formalism}.

\subsection*{Overlapping and Polyhierarchical Structure}

The models above assume each item belongs to exactly one node per
level. That is convenient for fitting and false of real query corpora:
a question about the thermodynamics of a chemical reaction is
legitimately a member of a physics-oriented and a chemistry-oriented
sub-domain at once. Soft or overlapping hierarchical clustering relaxes
this through a fuzzy or probabilistic membership
function~\parencite{fuzzyhc1969}; applied to a directed hierarchy the
relaxation yields a \textit{polyhierarchy}, in which a node may have
more than one parent and the structure is a directed acyclic graph
rather than a tree~\parencite{polyhierarchy2006}. Polyhierarchies are
standard in knowledge organisation and sparse in corpus-induced,
embedding-driven construction, where the DAG structure is typically
imposed post hoc rather than learned from overlap evidence.

\textbf{These are two separable claims, and this thesis retains the
first and rejects the second.}
Query assignment during fitting is weighted multi-membership under a
fixed posterior-responsibility floor, so a query holds non-negligible
membership in a small admitted set of leaves rather than diffuse
membership everywhere, with a primary (argmax) assignment retained as a
bookkeeping convention. The \emph{topology} stays a strict tree: every
node has exactly one parent, and overlap is carried by the membership
weights rather than by the edge set. The thermodynamics query therefore
reaches both leaves while neither leaf acquires a second parent.
A polyhierarchical variant was implemented and measured;
Section~\ref{sec:poly-negative} reports the measurement that retired it
and the cost of removing it.

\bigskip
\noindent
Three requirements for a query-level, unsupervised construction follow
from the above and are what Chapter~\ref{ch:system-architecture} builds
against: a node model that statistically characterises cluster geometry
(vMF); a splitting criterion with an explicit, chance-corrected
acceptance statistic rather than a heuristic threshold; and an
embedding space in which coarse and fine distinctions are both
resolvable (MRL).


\section{LLM-as-a-Judge and Reliability}\label{sec:llm-as-a-judge}

As human evaluation becomes a scalability bottleneck, LLM-as-a-judge
has become the dominant approach to automated assessment of model
outputs~\parencite{zheng2023judging}.
It uses the reasoning capabilities of frontier models to evaluate the
semantic quality of candidate responses. Its reliability properties,
and the conditions under which they hold, are therefore directly
relevant to any evaluation framework that delegates comparison to an
automated judge.

\subsection*{Pointwise and Pairwise Judging}

Pointwise evaluation asks a judge model to assign an absolute score to a
single response; pairwise evaluation presents two responses simultaneously
and asks for a preference judgment~\parencite{zheng2023judging}.
Zheng et al.\ introduced MT-Bench and Chatbot Arena to benchmark the two
framings against human preference, finding that pairwise evaluation
achieves higher agreement with human raters and is more robust to
prompt-formatting variation.
The relative nature of the task reduces the cognitive load of score
calibration, and that is the primary empirical justification for
preferring pairwise arena evaluation over pointwise scoring in
automated benchmarks.

\subsection*{Systematic Biases}

Three systematic biases distort pairwise preferences independently of
response quality.

\begin{itemize}
    \item \textbf{Positional bias.}
    Judges tend to favour the response presented first in the prompt,
    independently of content~\parencite{zheng2023judging}.
    Zheng et al.\ report a positional consistency of 65.0\% for
    GPT-4: in only about two thirds of pairs did the judge reach the
    same verdict when the positions of the two responses were swapped,
    implying that up to roughly 35\% of single-order pairwise outcomes
    can be position-driven rather than quality-driven.
    The magnitude strengthens rather than weakens the case for the
    bidirectional order evaluation adopted in this thesis: a bias that
    can touch a third of verdicts is a first-order term to be measured
    out, not a residual to be tolerated.

    \item \textbf{Verbosity bias.}
    Models consistently assign higher scores to longer responses, even
    when the additional content adds no reasoning
    depth~\parencite{chen2024verbosity}.
    This confound is particularly acute in knowledge-intensive QA,
    where a model can pad a correct but terse answer with
    confident-sounding elaboration to win a preference judgment it
    should not have won.

    \item \textbf{Self-preference bias.}
    LLMs exhibit a tendency to favour outputs stylistically similar to
    their own training distribution~\parencite{wataoka2025selfpreference,
    panickssery2024llm}, and \textcite{panickssery2024llm} show the
    effect is linked to self-recognition: evaluators that can identify
    their own generations are the ones that favour them.
    Recent evidence characterises this as a familiarity bias: a model
    assigns lower perplexity to outputs close to its own distribution and
    translates that fluency signal into a preference judgment, effectively
    rewarding self-similarity rather than correctness.
\end{itemize}

Mitigation strategies include position-swapping with consistency checks,
length-normalised scoring, and group-based polling to average out
idiosyncratic preferences~\parencite{zheng2023judging}.
Positional bias is the most tractable of the three, since a single
swap-and-consistency-check protocol identifies and discards
position-driven verdicts at low cost.
Verbosity bias is partially tractable through domain-specific prompt
engineering that makes evaluation criteria explicit, though the
effectiveness of this mitigation is an empirical question in any
given domain.
Self-preference bias remains the least tractable. It is structural
rather than prompt-addressable, and its magnitude varies with the
distance between the judging model's training distribution and
those of the evaluated models.

\subsection*{Judges on Verifiable Tasks: The Shortcut Result}

The biases above concern open-ended comparison. On tasks with a
verifiable answer, a further and more consequential behaviour is
documented, and it is the closest prior work to
Chapter~\ref{ch:results}'s central measurement.
\textcite{stephan2024calculation}, examining LLM judges on
mathematical reasoning tasks, find that judges tend to award the win
to the generally stronger model even on instances where that model's
answer is wrong; that a large fraction of judgments (70--75\%) can be
predicted from shallow features without consulting the mathematics at
all; and that judges are able to detect the stratum in which both
candidate answers are correct --- the stratum where a preference
verdict is least informative. On verifiable tasks, in short, verdicts
track answer-level correctness and prior model strength rather than
the quality of the reasoning being adjudicated.

Meta-evaluation benchmarks corroborate the pattern from the other
direction. \textit{JudgeBench}~\parencite{tan2025judgebench}
constructs response pairs whose preference label is objective
correctness and finds that even strong judge models perform near
chance on them; \textit{LLMBar}~\parencite{zeng2024llmbar} shows
judges systematically fail pairs in which superficial quality and
instruction adherence are dissociated; and
\textit{RewardBench}~\parencite{lambert2024rewardbench} standardises
this style of adversarial audit for reward models.
Judge competence, these benchmarks agree, is not a scalar that
transfers from fluent domains to verifiable ones.

Chapter~\ref{ch:results}'s central finding is therefore positioned as
a pre-registered \emph{extension} of this line, not an independent
discovery: relative to \textcite{stephan2024calculation} it moves the
design from free-form mathematics to multiple choice, where the answer
key is exact rather than extracted; it tests the rubric mechanism
directly, through correctness strata the rubric cannot distinguish;
and the strata analysis was registered before any verdict existed
(Chapter~\ref{ch:experimental-design}).

\subsection*{Domain Calibration}

A judge prompted with a generic instruction to ``prefer the better
answer'' applies implicit quality criteria drawn from its training
distribution, which may not align with the domain-specific standards
relevant to the query.
\textit{FLASK}~\parencite{ye2024flask} makes this concrete:
decomposing coarse quality judgments into fine-grained,
instance-specific skill criteria changes both the scores and the
resulting model orderings, so the choice of criteria is part of the
measurement rather than a presentational detail.
This misalignment is not merely a calibration error in the statistical
sense; it changes the semantic content of the preference judgment.
In mathematical reasoning, for instance, a judge without domain
conditioning will prefer a fluent but incorrect derivation over a
terse but correct one, conflating presentation quality with
mathematical validity.

The earliest mitigation in this family is Zheng et al.'s
reference-guided mode, in which the judge is shown a reference
solution at verdict time~\parencite{zheng2023judging}.
In industrial practice, \textit{MemAlign} --- a non-peer-reviewed
system description~\parencite{databricks2026memalign}, cited here as
an industrial instance only --- aligns synthetic judges with
domain-specific standards through a dual-memory design (general
evaluation principles plus episodic edge cases) calibrated against as
few as 30--50 human gold labels without weight updates.
On the training side, the companion RL-trained reasoning judges
\textit{Think-J}~\parencite{huang2025thinkj} and
\textit{JudgeLRM}~\parencite{chen2025judgelrm} reward explicit
chain-of-thought justifications that are consistent with the final
preference, making the evaluation process auditable and reducing the
proportion of preference reversals attributable to surface features.
These approaches indicate that domain-specific judge conditioning
improves alignment with human expert ratings, at the cost of requiring
either labelled examples or a trained reasoning judge.

The setting that arises in TaxoArena sits between these mitigations
and shares none of their supervision.
Its judge-generation phase has access to MMLU-Pro reference answers
from the construction set but not to human preference annotations,
while the arena evaluation phase has access to neither.
This differs from Zheng et al.'s reference-guided mode, which reveals
the reference at verdict time, and from FLASK, whose criteria are
instantiated per instance at evaluation time: here a rubric is
distilled \emph{once} from the answer key at construction, and the key
is withheld at test.
The narrow claim this thesis makes is not that reference-informed
judging is unstudied; it is that this distill-once,
withhold-at-test protocol had not been tested against correctness
strata, which is the design under which Chapter~\ref{ch:results}
measures it.

\subsection*{The Question This Thesis Measures}

The biases above are all cited results, and this thesis claims novelty
for none of them. That LLM judges prefer longer answers, favour the
first-presented answer, and fall back on undocumented heuristics where
a rubric is silent is established --- and so, on verifiable tasks, is
the shortcut to answer correctness and prior model
strength~\parencite{stephan2024calculation}.

What Chapter~\ref{ch:results} measures is the question this shortcut
literature opened, moved to the setting where it is sharpest:
\emph{on a corpus with a verifiable answer key, what does a pairwise
judge condition its verdict on?} The domain-calibration literature
assumes, reasonably, that a better-specified rubric changes the
judgment. That assumption is testable whenever the corpus supplies an
independent correctness signal, because a judge that is really applying
criteria and one that is verifying an answer will behave differently
under strata the criteria cannot distinguish --- items where both
responses are correct, items where neither is, items where neither
response contains any reasoning to evaluate. Those strata exist in
MMLU-Pro, cost nothing to construct, and --- unlike the post-hoc
predictability analyses of prior work --- were registered here before
the verdicts existed. The pre-registered stratified decomposition
reported in Section~\ref{sec:res-judge} is the contribution; the
shortcut phenomenon itself is prior art, and the individual biases are
background.


\bigskip
\noindent
The three biases and the domain-calibration requirement constrain the
judge design in any pairwise arena framework.
Positional bias requires presentation-order randomisation and
consistency filtering; verbosity bias requires explicit per-domain
criteria that do not reward length, whose effectiveness is an empirical
question; self-preference bias remains a residual structural limitation
unless judge models are rotated.
The guideline-generation mechanism that follows from these constraints
is specified in Chapter~\ref{ch:system-architecture}.

Such a judge produces one verdict per comparison, with ties represented
as half-wins to each model, following the tie-as-half-win convention
introduced for paired-comparison models by
Davidson~\parencite{davidson1970ties}.
These are the observations consumed by the ranking models reviewed in
the following section.



\section{Pairwise Ranking Models}\label{sec:ranking-models}

This section reviews the statistical foundations of pairwise ranking,
establishes the limitations of the dominant scalar approach, and
contextualises active scheduling within the theoretical literature on
sample-efficient ranking.

\subsection*{Elo and Its Limitations}

The \textit{Elo rating system}~\parencite{elo1978rating} is the
starting point for comparative evaluation.
Developed for zero-sum games like chess, it assumes that an agent's performance
is a random variable centred around a latent skill level $R$.
The probability $E_A$ of agent $A$ defeating agent $B$ follows a logistic
distribution~\parencite{boubdir2024elo}:
\begin{equation}
    E_A = \frac{1}{1 + 10^{(R_B - R_A)/400}}
    \label{eq:elo-expected}
\end{equation}
While Elo is reliable for tracking stable skill levels over time, it suffers from
several weaknesses in LLM evaluation contexts.
It assumes static skill between matches and does not explicitly model uncertainty,
leading to ranking instability in populations with high variance or sparse
interaction data~\parencite{boubdir2024elo}.
Elo ratings are also transitive by construction, assuming that
if $A \succ B$ and $B \succ C$, then $A \succ C$.
\textcite{boubdir2024elo} show that LLMs frequently exhibit
non-transitive preferences, which renders linear Elo leaderboards
statistically misleading.

\subsection*{Bradley--Terry and Bayesian Extensions}

The \textit{Bradley--Terry} (BT) model provides the formal probabilistic
link for pairwise preferences, modelling the merit of each model as a
latent parameter $\lambda_i > 0$~\parencite{bradley1952rank}.
The probability that model $i$ is preferred over model $j$ is
\begin{equation}
    P(i \succ j) = \frac{\lambda_i}{\lambda_i + \lambda_j}
    = \sigma(\theta_i - \theta_j),
    \label{eq:bt}
\end{equation}
where $\theta_i = \log \lambda_i$ and $\sigma$ is the logistic function.
The log-likelihood is strictly concave in $(\theta_i)$, and the MLE
$\hat{\boldsymbol{\theta}}$ is obtained via iterative re-weighted least
squares (IRLS) or the MM algorithm~\parencite{bradley1952rank}.
The asymptotic covariance of $\hat{\boldsymbol{\theta}}$ is the inverse of
the Fisher information matrix, yielding per-parameter standard errors
$\hat{\sigma}_i$.
A necessary condition for identifiability is that the comparison
graph---the graph in which each node is a model and each edge
represents at least one observed comparison---is connected; if it decomposes
into two or more components, BT parameters are identified only up to an
additive constant per component and cross-component comparisons cannot
be inferred~\parencite{bradley1952rank}.

This identifiability requirement is non-trivial in per-domain arenas,
where the query pool at a leaf may be small and not all model pairs
are guaranteed to receive a comparison.
It nonetheless becomes a structural advantage of the per-domain
decomposition in TaxoArena. Restricting the comparison graph to the
small, semantically coherent query pool of a single leaf leaves the
bootstrap phase needing only $\binom{M}{2}$ initial comparisons, one
per pair, to guarantee connectivity, a number that is tractable even
for leaves with modest query populations.
Graph connectivity is necessary but not sufficient for a finite
unregularised MLE to exist, since complete separation (one model
winning every comparison against another) prevents convergence of
the MM algorithm; the regularisation strategy used to handle this
case is specified in Chapter~\ref{ch:methodology}.

Bayesian rating systems such as
\textit{TrueSkill}~\parencite{herbrich2007trueskill} and the
Weng--Lin approximations implemented in
OpenSkill~\parencite{weng2011bayesian} represent
skill as a Gaussian with explicit uncertainty and converge faster than
Elo by weighting each match by its information gain.
Their online-update design targets streaming match data; for a
fixed-budget batch arena, the BT MLE with Fisher standard errors
provides the same uncertainty quantification within a simpler and more
interpretable estimation framework.
A further extension is hierarchical, partial-pooling Bayesian BT
modelling, in which per-group strength parameters shrink toward a
population-level mean~\parencite{cattelan2012paired}; it is the
natural refinement for a per-domain arena whose smallest cells carry
few comparisons, and this thesis instead estimates each cell's BT
model independently, keeping per-cell profiles interpretable in
isolation at the cost of borrowing no strength across cells.

\subsection*{Active Ranking and Sample Complexity}

The na\"ive strategy for recovering a full ranking of $M$ models using pairwise
comparisons is round-robin evaluation, requiring $\binom{M}{2}$ comparison
pairs per leaf.
\textit{Active ranking} addresses this by selecting, at each step, the
comparison most likely to reduce global ranking uncertainty.

Jamieson and Nowak~\parencite{jamieson2011activeranking} established the
theoretical basis, showing that if $n$ items can
be embedded in a $d$-dimensional metric space consistent with the ranking, then
a sequentially adaptive algorithm recovers the full ranking using only
$O(d \log n)$ pairwise comparisons on average, an exponential reduction over
the $\Theta(n^2)$ required by non-adaptive strategies.
The same work proves that \textit{random} (non-adaptive) sampling requires
almost all $\binom{n}{2}$ comparisons to identify the ranking with
reasonable probability, establishing a fundamental separation between adaptive
and non-adaptive query selection.
The metric-space embedding assumption underlying this result---that model
capabilities can be represented as points in a low-dimensional Euclidean
space whose order is consistent with the true ranking---is an idealisation.
LLM capabilities are likely not perfectly embeddable in a fixed
low-dimensional metric space; restricting the ranking problem to a
semantically coherent sub-domain reduces the effective intrinsic
dimensionality and makes the metric-space assumption more defensible,
which is a requirement for the theoretical guarantee to be practically
meaningful.

In the stochastic setting, where comparison outcomes are noisy Bernoulli
random variables as in an LLM judge, Heckel et al.~\parencite{heckel2019activeranking}
derive finite-sample bounds under the Bradley--Terry model directly.
Their main result characterises the number of comparisons required to
distinguish two adjacent models in the ranking with error probability $\delta$:
the sample complexity scales as $O(\Delta^{-2} \log(M/\delta))$, where
$\Delta$ is the minimum separation between consecutive BT parameters.
Two consequences follow.
First, pairs whose BT parameters are close together require
disproportionately many comparisons; a uniform budget allocation is
therefore inefficient.
Second, a confidence-interval-based allocation rule that schedules
comparisons for the pair with the largest overlap in BT confidence
intervals is near-optimal, achieving the
$O(\Delta^{-2} \log(M/\delta))$ bound without knowledge of
$\Delta$~\parencite{heckel2019activeranking}.
Heckel et al.\ further establish that imposing a parametric model (BTL,
Thurstone) offers at most $O(\log n)$ improvement over non-parametric
adaptive strategies for stochastic comparisons, indicating that the
primary practical advantage of the BT model lies in its interpretability
and its natural provision of a convergence criterion via standard errors,
rather than in any asymptotic sample-complexity gain.

\subsection*{Multidimensional Ranking and Intransitivity}

Cyclic dominance, the ``Rock-Paper-Scissors'' effect, arises when models
have specialised capabilities that excel in certain domains but fail in
others, violating the transitivity assumption of scalar
models~\parencite{balduzzi2018melo}.
To model these interactions, current research uses vectorised approaches.
\textit{mELO}~\parencite{balduzzi2018melo} applies Hodge Decomposition to
separate model interactions into a transitive (scalar) component and a cyclic
(skew-symmetric) component, allowing more accurate win-rate predictions where
models exhibit diverse, niche strategies.
The \textit{Blade-Chest model}~\parencite{chen2016bladechest} infers a
$d$-dimensional representation for each model, modelling ``offence'' and
``defence'' separately; by learning a dataset-specific metric space, it can
represent any pairwise stochastic preference relation.
Both lines of work target the same failure of scalar ratings, that a
single global ranking cannot represent models which excel in some
domains and fail in others, and both report more accurate pairwise
outcome predictions once interactions are modelled directly.

\bigskip
\noindent
Taken together, the ranking literature identifies three requirements for
an arena evaluation framework: a probabilistic pairwise model with explicit
uncertainty quantification and a connected comparison graph as a
precondition for identifiability; an adaptive scheduler that concentrates
comparisons on the hardest-to-distinguish pairs, achieving near-optimal
sample complexity without knowledge of the true parameter separations; and
a per-domain decomposition that makes the transitivity assumption locally
defensible and reduces the effective dimensionality required by the
Jamieson--Nowak metric-space bound.
These requirements directly motivate the routing use case examined in
Section~\ref{sec:routing}.


\section{Query Routing as the Target Use Case}\label{sec:routing}

Query routing is the downstream task that per-domain capability
profiles are designed to enable. Routing systems that induce their own
domain structure from a query corpus now exist, so the gap motivating
this thesis is not the absence of induced structure; it is the absence
of a routing substrate that is simultaneously interpretable,
statistically characterised, and validated independently of routing
accuracy. This section names the occupying systems and states that
narrower gap.

\subsection*{The Routing Problem in Agentic Systems}

In multi-model agentic systems, a query is dispatched to one of several
candidate models based on an estimate of which model will perform best
on that query.
The optimal dispatch policy requires knowledge of each model's
domain-conditional performance. If model $A$ outperforms model $B$ on
mathematical reasoning but underperforms on legal analysis, a global
ranking that averages across domains will systematically mis-route
queries at both extremes.
This is the case for \textit{domain-centric
evaluation}~\parencite{raju2024constructing}. Raju et al.\ show on
chat-assistant evaluation that model rankings are not stable across
domains, and that a router informed by per-domain scores substantially
reduces response-quality variance relative to global-rank routing.
This line of work, however, requires that per-domain scores already
exist and that the domain taxonomy is either given or constructed by
hand, which limits reproducibility and scalability.

\subsection*{Learned Routing Systems}

\textit{FrugalGPT}~\parencite{chen2023frugalgpt} routes by cascading
through models of increasing capability until a learned confidence
estimator accepts the current response, reducing inference cost by up
to 98\% at matched accuracy.
\textit{RouteLLM}~\parencite{ong2024routellm} instead trains binary
routing functions, among them a similarity-weighted Bradley--Terry
router, on 80K pairwise preference judgements from Chatbot Arena.
Routers trained solely on the global preference
distribution generalise poorly to domain-specific benchmarks such as
MMLU and GSM8K unless the training data is augmented with in-domain
preference labels~\parencite{ong2024routellm}, because a globally
trained BT model does not encode the domain-conditional capability
variation that accurate routing requires.
\textit{EmbedLLM}~\parencite{zhuang2024embedllm} and
\textit{UniRoute}~\parencite{jitkrittum2025uniroute} route
geometrically, representing each model by an embedding derived from
its performance on a representative query set; the quality of that
per-model profile bounds achievable routing accuracy.
In FrugalGPT, RouteLLM, and EmbedLLM, the domain structure over which
models are profiled is either absent (a global ranking) or predefined
by the benchmark used to construct the training set.

A more recent line induces its routing structure directly from the
query corpus, and it forecloses any broad novelty claim here.
\textit{The Avengers}~\parencite{zhang2025avengers} embeds and
clusters training queries, scores each candidate model per cluster,
and routes a test query to the models that dominate its nearest
clusters; \textit{Avengers-Pro}~\parencite{zhang2025avengerspro}
extends the recipe with a tunable performance--efficiency trade-off.
UniRoute's cluster-based variant represents each model by its error
profile over query clusters for exactly the same
reason~\parencite{jitkrittum2025uniroute}, and
\textit{RouterDC}~\parencite{chen2024routerdc} learns query and model
embeddings jointly by dual contrastive learning, inducing a soft
structure in the shared space. Corpus-induced structure for routing is
therefore occupied territory; these systems are the competing
baselines for TaxoArena's routing use case, and standardised router
evaluations such as \textit{RouterBench}~\parencite{hu2024routerbench}
make that comparison feasible.

What none of these systems provides --- and what P2L
(Section~\ref{sec:benchmarks}), the strongest occupant of the
per-domain-profile territory, provides only in continuous, amortised
form~\parencite{frick2025p2l} --- is the conjunction this thesis
targets: routing structure that is (a) discrete and
human-interpretable, seeded from editorial labels and then refined by
geometric evidence; (b) statistically characterised per cell, with a
tested coherence model on the query side and Bradley--Terry profiles
carrying Fisher-information confidence intervals on the model side;
and (c) built as an \emph{instrument}, its cells validated for
coherence independently of --- and prior to --- any routing
objective, an instrument/router separation that cluster-based
routers, which tune their structure against routing accuracy itself,
do not have.
The Avengers, Avengers-Pro, UniRoute's cluster mode, and RouterDC
occupy the structure-induction territory; the narrowed gap is the
conjunction, not the induction.

\section*{Synthesis}

The five bodies of literature reviewed in this chapter converge on a
single compound problem. Corpus-induced taxonomies for evaluation
exist (BenchBuilder, TnT-LLM, Clio), corpus-induced routing structures
exist (The Avengers, UniRoute's cluster mode, RouterDC), and
per-prompt capability profiles with uncertainty exist (P2L); no
existing system, however, combines a generative, coherence-tested node
model over a discrete, human-interpretable partition with per-cell
Bradley--Terry profiles carrying formal confidence intervals,
validated as an instrument before any router consumes them.
Both the ranking literature and the routing literature require exactly
that conjunction as substrate.
TaxoArena is built to test whether that substrate can exist: a
label-bootstrapped, corpus-induced partition supplies the cells, a
reference-informed, answer-key-blind arena judges inside them, and ---
because the corpus carries a verifiable key --- every step of the
chain from coherent cells to better judgments is separately
measurable. Chapter~\ref{ch:results} reports which links held.
The system design is specified in Chapter~\ref{ch:system-architecture};
the formal foundations are derived in Chapter~\ref{ch:methodology}.
